\section{Introduction}

Finger-intent prediction from surface electromyography (sEMG) is an appealing interface for post-stroke neurorehabilitation because it offers a non-invasive path from residual muscle activity to assistive control, therapy feedback, and fine-motor retraining~\citep{munoznovoa2022upper,zheng2021semg,boukhennoufa2022wearable}. In the broader PSR project, that decoding capability is intended to support machine-learning models that can eventually be served on Raspberry Pi 5-class hardware and coupled to a rehabilitation glove. This manuscript focuses on the software research layer of that agenda: the data handling, signal processing, modeling, and evaluation work that turns multichannel sEMG into five-finger intent predictions.

The challenge is not only raw predictive accuracy, but also methodological reproducibility. Post-stroke sEMG is noisy, patient-specific, and highly sensitive to recording conditions, impairment severity, and dataset construction, so benchmark results depend heavily on how signals are filtered, denoised, segmented, labeled, tensorized, and split across patients. Reviews of sEMG-based gesture recognition consistently note that preprocessing design, feature representation, and deployment constraints can matter as much as the final classifier family~\citep{li2021gesture,ni2024survey}. A research project that reports model performance without a stable preprocessing and evaluation path is therefore difficult to audit and even harder to extend toward embedded rehabilitation use.

\repo\ addresses that gap by providing one auditable software framework for the project~\citep{psrrepo2026}. The current implementation includes \physiomio\ ingestion, gesture-to-finger label mapping, a fixed signal-conditioning stack, window-level feature extraction, shared tensor adapters, direct LSTM/CNN/GNN baselines, a merged 1D CNN student/teacher suite, Optuna-based hyperparameter optimization, healthy-to-impaired transfer-learning summaries, and benchmark artifacts stored as committed metric files. Recent project updates also added teacher-ensemble distillation, per-finger threshold tuning, and latency benchmarking utilities, which materially strengthen the edge-deployment research story even when the outputs of those utilities are not yet fully committed as benchmark bundles.

The central story of this paper is therefore research-oriented rather than repository-oriented. We do not present a new decoding architecture, nor do we claim that the current implementation already demonstrates end-to-end rehabilitation hardware validation. Instead, we report a software research program whose current contribution lies in making experimental assumptions explicit: which recordings are used, how they are preprocessed, how tensors are shared across model families, how baseline and improved architectures are compared, which metrics are computed, and which deployment-facing hooks are already present in code.

This paper makes four study-level contributions. First, it formalizes the current software pipeline from \physiomio\ ingestion through preprocessing, feature extraction, tensor adaptation, model fitting, and deployment-screening utilities. Second, it establishes a clear benchmark narrative that begins with three direct model families, LSTM, CNN, and GNN, under one preprocessing regime before analyzing subsequent improvements. Third, it evaluates two concrete improvement pathways now visible in the repository: compact CNN student and ResNet teacher development, and healthy-to-impaired transfer learning, while also documenting distillation as an implemented but not yet fully benchmarked speed-quality path. Fourth, it situates the project against the broader literature on post-stroke sEMG decoding, including PhysioMio, Ninapro, and recent transfer-learning studies, while drawing a careful boundary around what the current evidence still cannot support, especially distilled-student claims, on-device Raspberry Pi 5 validation, and clinical outcome conclusions.
